%% AMS LaTeX Package Version 6.2 (ametsocV6.2.cls, ametsocV6.bst).
%% Target: Artificial Intelligence for the Earth Systems (AIES), submission
%% type "Lessons Learned": short papers on insights into how well AI methods
%% work for a whole class of Earth-system applications.
%%
%% Ported from the arxiv.sty skeleton (kept in _trash/main_arxiv_2026-11-09.tex)
%% on 2026-11-09. The class has no journal option and no keywords command;
%% keywords are entered in the submission system.
%%
%% Conventions:
%%   \todo{}  text or number still to write or measure (red)
%%   \ev{}    where the number comes from (blue). Paths are relative to
%%            rainshift-uda/ or results_rainshift_uda/. Delete before
%%            submission: \renewcommand{\ev}[1]{} below hides them all.
%% No result numbers are carried over from the old project description:
%% most were withdrawn (NOTES.md, "Graveyard").
%%
%% Keywords (for the submission system): domain adaptation; precipitation
%% downscaling; super-resolution; distribution shift; evaluation.

%%%%%%%%%%%%%%%%%%%%%%%%%%%%%%%%%%%%%%%%%%%%%%%%%%%%%%%%%%%%%%%%%%%%%
% PREAMBLE
%%%%%%%%%%%%%%%%%%%%%%%%%%%%%%%%%%%%%%%%%%%%%%%%%%%%%%%%%%%%%%%%%%%%%

% 1.5-spaced peer-review format (the only format AMS accepts for submission)
\documentclass{ametsocV6.2}
% Two-column journal layout, for our own reading only:
% \documentclass[twocol]{ametsocV6.2}

% The class already loads graphicx, amsmath/amssymb, natbib, url, xcolor.
% Tables use the class rules \topline, \midline, \botline (no booktabs).
\usepackage{tabularx}

\newcommand{\todo}[1]{\textcolor{red}{[TODO: #1]}}
\newcommand{\ev}[1]{\textcolor{blue}{[Evidence: #1]}}
% \renewcommand{\ev}[1]{}   % uncomment to hide evidence pointers
\newcommand{\lamhat}{\hat{\lambda}}

%%%%%%%%%%%%%%%%%%%%%%%%%%%%%%%%%%%%%%%%%%%%%%%%%%%%%%%%%%%%%%%%%%%%%

\title{What unlabelled target data cannot fix: lessons from domain\\
adaptation for cross-region precipitation downscaling}
% Alternative titles:
%  - Lessons learned from unsupervised domain adaptation for precipitation
%    super-resolution across regions
%  - Before adapting, measure: lessons from domain adaptation for
%    precipitation downscaling

\authors{Filippo Quarenghi\aff{a}\correspondingauthor{Filippo Quarenghi, filippo.quarenghi@unil.ch}
and Tom Beucler\aff{a}}

\affiliation{\aff{a}{Faculty of Geosciences and Environment and Expertise Center for
Climate Extremes, University of Lausanne, Lausanne, Switzerland}}

%%%%%%%%%%%%%%%%%%%%%%%%%%%%%%%%%%%%%%%%%%%%%%%%%%%%%%%%%%%%%%%%%%%%%
% ABSTRACT (no more than 250 words; no citations)
% Shape: problem -> what we did -> the lessons -> the takeaway.
%%%%%%%%%%%%%%%%%%%%%%%%%%%%%%%%%%%%%%%%%%%%%%%%%%%%%%%%%%%%%%%%%%%%%

\abstract{Neural downscaling of precipitation is trained where
high-resolution observations exist and applied where they do not.
Unsupervised domain adaptation (UDA) promises to close the resulting
transfer gap with unlabelled predictors from the target region. We report
what we learned testing that promise on \todo{N} UDA methods, a
deterministic U-Net and a flow-matching generative model, on three
RainShift regions (ERA5 to IMERG) and on a perfect-model COSMO-CLM
setting. \todo{One sentence per lesson, five sentences, written last.}
The common thread is that UDA can only help when the
precipitation-given-atmosphere relation is shared between regions, and
that this, together with what labelled data could achieve, can and should
be measured before any adaptation is trained.}

\begin{document}

%% Necessary!
\maketitle

%%%%%%%%%%%%%%%%%%%%%%%%%%%%%%%%%%%%%%%%%%%%%%%%%%%%%%%%%%%%%%%%%%%%%
% SIGNIFICANCE STATEMENT (no more than 120 words, plain language)
%%%%%%%%%%%%%%%%%%%%%%%%%%%%%%%%%%%%%%%%%%%%%%%%%%%%%%%%%%%%%%%%%%%%%

\sigstatement{Machine-learning models that sharpen coarse rainfall
fields are trained where dense observations exist, mostly in Europe and
North America, but are needed most where observations are sparse. A
family of techniques claims to adapt such models to a new region using
only its weather data, without its rainfall observations. We find that
these techniques rarely help, and we explain when they cannot: if the
same weather produces different rainfall in the two regions, no amount of
unlabelled weather data reveals the difference. We give practitioners
checks to run before investing in adaptation, and list the evaluation and
implementation pitfalls that made failing methods look successful in our
own experiments.}

%%%%%%%%%%%%%%%%%%%%%%%%%%%%%%%%%%%%%%%%%%%%%%%%%%%%%%%%%%%%%%%%%%%%%
% MAIN BODY
%%%%%%%%%%%%%%%%%%%%%%%%%%%%%%%%%%%%%%%%%%%%%%%%%%%%%%%%%%%%%%%%%%%%%

\section{Introduction}\label{sec:intro}
% Target: ~1 page. Four paragraphs and the lesson list.

% P1 -- The deployment problem.
High-resolution precipitation observations are dense over Europe, North
America and East Asia and sparse over much of the tropics, where
downscaled information is most needed. Neural downscaling models lose
much of their skill when moved between regions
\citep{harder2025rainshift}. \todo{cite Harder et al. arXiv:2512.01400
(WGAN, 15 areas).} Unsupervised domain adaptation (UDA) offers an
apparently cheap remedy: keep the labelled source region, add unlabelled
predictors from the target region, and align the two during training
\citep{ganin2016domain,long2015learning,sun2016deep}.

% P2 -- Why it might not work: theory in one paragraph.
The standard justification bounds the target risk by the source risk, a
divergence between input distributions, and the risk $\lambda$ of the best
hypothesis on both domains
\citep{bendavid2010theory,redko2017theoretical}. Alignment acts only on
the divergence. It can help under covariate shift, where $P(X)$ changes
and $P(Y\mid X)$ does not. When $P(Y\mid X)$ itself differs, the target
conditional is not identifiable from unlabelled target inputs, and
enforcing invariance can raise target error
\citep{zhao2019invariant,johansson2019support}. \todo{cite Kong et al.
2023, arXiv:2306.06510, for non-identifiability of the joint.}

% P3 -- The literature disagrees, and why (the organising thesis).
Reports on cross-region transfer in downscaling disagree. We argue that
they sort by one property of the experimental design: whether the
conditional is shared by construction. Perfect-model and same-product
designs, in which the coarse input is derived from the high-resolution
target, share it, and transfer works \citep{qian2026}. \todo{cite Singh
et al. 2026 GMD; CORDEX-ML-Bench arXiv:2606.29172; our EGU26-19396
abstract.} Reanalysis-to-observation designs such as RainShift do not
share it by construction, and transfer fails
\citep{harder2025rainshift}. We test both kinds of design with the same
models and methods.

% P4 -- What this paper is, and the lessons.
This is a lessons-learned paper, not a method paper. We trained \todo{N}
models across \todo{N} transfer directions, two architectures and
\todo{N} UDA methods, and audited our own pipeline when results looked
too good. The lessons, developed in section~\ref{sec:lessons} and
collected as a checklist in Table~\ref{tab:checklist}, are:
\begin{enumerate}
  \item Measure whether the conditional shifts before choosing to adapt,
        and calibrate that measurement against a null that keeps the
        data's dependence structure (section~\ref{sec:L1}).
  \item Bound the problem with labels first: a jointly trained oracle and
        a few-shot curve say what any adaptation could gain, and at what
        price in labels (section~\ref{sec:L2}).
  \item Check that the adaptation term acts as intended; inert losses,
        mixed normalisation statistics and sign errors are silent
        (section~\ref{sec:L3}).
  \item Evaluation choices can manufacture recovery: report trivial
        references, physical-unit diagnostics and seed spread beside every
        transformed-space score (section~\ref{sec:L4}).
  \item A probabilistic objective does not change the diagnosis, but it
        changes what the score rewards (section~\ref{sec:L5}).
\end{enumerate}

%%%%%%%%%%%%%%%%%%%%%%%%%%%%%%%%%%%%%%%%%%%%%%%%%%%%%%%%%%%%%%%%%%%%%
\section{Setting}\label{sec:setting}
% Target: ~1.5 pages. Only what the lessons need; details to appendix A.

\subsection{Two transfer problems}\label{sec:data}

\paragraph{RainShift: reanalysis to observations.} RainShift
\citep{harder2025rainshift} pairs ERA5 predictors
\citep{hersbach2020era5} with IMERG precipitation
\citep{huffman2023imerg}. We use three regions spanning a difficulty
axis: western Europe, the Horn of Africa and Melanesia, giving six
ordered transfer directions. \todo{inputs (9 dynamic + 2 static), grid
sizes, time resolution and periods, split sizes; wind level.}
\ev{overview\_2026-11-07.md \S3; memory split-sizes}

\paragraph{COSMO-CLM: a perfect model.} \todo{COSMO-CLM 2.2\,km
(Cui et al. 2024); coarse inputs by 10$\times$ block mean of the
high-resolution fields, stated plainly; five European regions (southern
Scandinavia, central Spain, Alps, British Isles, Poland).} One model, one
product and no observation network: the conditional can still move
between regions through the physics, but none of RainShift's product
confounds remain. \ev{PLAN\_COSMO.md Part B, C0--C8. C1 built 2026-11-09;
C2 onward not run.}

\paragraph{Normalisation.} Every dataset, target included, is
standardised with the \emph{source} statistics. Normalising the target
with its own statistics would remove part of the covariate shift being
studied. \todo{one sentence on the log transform.}

\subsection{Models and adaptation methods}\label{sec:methods}

\paragraph{Models.} A dual-encoder U-Net trained with mean squared error
in log-precipitation space, and adaptive flow matching (AFM)
\citep{fotiadis2025adaptive,lipman2023flow} with the same encoder, trained
end to end and scored with the fair CRPS
\citep{gneiting2007scoring,ferro2014fair}. \todo{widths, batch, LR,
selection rule; why width 32.}

\paragraph{UDA methods.} Table~\ref{tab:methods} groups the methods by
where they act and the assumption under which they can help.

\begin{table}[t]
\caption{Methods, where they act, and the assumption under which they can
reduce target risk. CORAL and multi-scale MMD were dropped after audit
(appendix B).}\label{tab:methods}
\begin{center}
\begin{tabularx}{\linewidth}{llX}
\topline
Acts on & Method & Can help if \\
\midline
Inputs    & FDA \citep{yang2020fda}; quantile mapping (RainShift baseline)
          & $P(Y\mid X)$ shared (covariate shift) \\
Features  & DANN \citep{ganin2016domain}; MMD \citep{gretton2012kernel,long2015learning}
          & a representation exists in which $P(Y\mid\phi(X))$ is shared \\
Joint     & JDOT \citep{courty2017jdot,damodaran2018deepjdot}
          & the target conditional is close to a transported source one \\
Outputs   & spectral (PSD) matching
          & spatial structure is shared \\
Test time & AdaBN \citep{li2018adabn}
          & the shift is captured by first and second moments of activations \\
\midline
Labels (reference) & oracle (joint training); few-shot fine-tuning & --- \\
\botline
\end{tabularx}
\end{center}
\end{table}

\subsection{Protocol and decomposition}\label{sec:protocol}

For each direction $S\to T$ we train a source-only model $f_S$, the
target's in-domain model $f_T$, an oracle $h_{ST}$ trained on both label
sets, and the UDA variants, each with three seeds. On the target test
set,
\begin{equation}
\underbrace{\mathcal{R}_T(f_S)-\mathcal{R}_T(f_T)}_{\text{transfer gap}}
=
\underbrace{\mathcal{R}_T(f_S)-\mathcal{R}_T(h_{ST})}_{\text{budget}}
+
\underbrace{\mathcal{R}_T(h_{ST})-\mathcal{R}_T(f_T)}_{\lamhat},
\label{eq:decomp}
\end{equation}
where $\lamhat$ is a class-restricted estimate of the part of the gap no
single model in the class removes, and the budget is what adaptation
could recover. A method's recovery is the fraction of the budget it
closes; its source cost is the change in $\mathcal{R}_S$.
\todo{state that seeds vary initialisation only (the dataset seed is
fixed), and the two hardware/software stacks.}

\paragraph{Metrics.} The headline score is the squared error of
$\log(P+0.1\,\mathrm{mm\,h^{-1}})$ for the U-Net and the CRPS in the same
space for AFM. Every such score is reported beside the mm-space bias, the
wet-pixel fraction and the field maximum (lesson~4).
\ev{D1, 2026-10-29; floor literature in overview \S3}

%%%%%%%%%%%%%%%%%%%%%%%%%%%%%%%%%%%%%%%%%%%%%%%%%%%%%%%%%%%%%%%%%%%%%
\section{Lessons}\label{sec:lessons}
% Each lesson: (i) the lesson in one sentence, (ii) what happened to us,
% (iii) the evidence (one table or figure at most), (iv) the
% recommendation, generalised beyond this task. Target: ~1 page each.

%---------------------------------------------------------------------
\subsection{Lesson 1: Measure the conditional shift before adapting}
\label{sec:L1}

\paragraph{Lesson.} Whether UDA has anything to do depends on whether
$P(Y\mid X)$ differs between regions. That can be tested directly from
labelled data in both regions, before any adaptation is trained, but only
against a null that respects the data's spatial and temporal dependence.

\paragraph{What happened.} \todo{Matched-input test on (CAPE, TCW) bins;
our first pixel-permutation null was 5--10$\times$ too small and inflated
the excess factors, which we withdrew. Conditional-permutation null with
time blocks: Europe--Horn and Europe--Melanesia differ in the
conditional; Horn--Melanesia is undetermined, and at limited support the
conditional and support effects cannot be separated. A model-based 9-D
test had no power for lack of overlap.}
\ev{P1, P1b: unet/conditional\_shift\_binned\_p1b.json; P2:
unet/conditional\_shift\_model\_p2/; overview \S2 claim 1, \S5.4}
\todo{COSMO C2: the same test across the COSMO region pairs.}

\paragraph{Recommendation.} \todo{Use a conditional test with a
dependence-preserving null and positive controls; report common support
next to it; treat ``undetermined'' as an answer.}

% Figure 1 candidate: matched-input statistic against its null, RainShift
% pairs and COSMO pairs side by side.

%---------------------------------------------------------------------
\subsection{Lesson 2: Bound the problem with labels first}\label{sec:L2}

\paragraph{Lesson.} The oracle of Eq.~(\ref{eq:decomp}) and a few-shot
curve are cheap compared with a method sweep, and together they say how
much any adaptation could gain and how many target labels would buy the
same gain.

\paragraph{What happened.} \todo{Oracle: $\lamhat$ small into the tropics,
22--39\,\% of the gap into western Europe for the U-Net, mostly gone for
AFM. Few-shot: 2000 labelled target samples reach the oracle in every
direction. So the gap is reducible in principle and cheaply closed with
labels, yet not identifiable without them.}
\ev{P14 table: results/seed<S>/field\_mean\_check\_qm/; oracle trees
unet\_seed<S>\_oracle/, afm\_seed42\_oracle/; P16:
results/seed<S>/unet\_fewshot/; overview \S5.1, \S5.6}
\todo{COSMO C4--C5.}

\paragraph{Recommendation.} \todo{Before a UDA study, run the oracle and
a few-shot curve; if a few hundred labels close the gap, the practical
question is label acquisition, not adaptation.}

% Figure 2 candidate: per direction, the gap split into budget and
% lambda-hat, with the few-shot points overlaid.

%---------------------------------------------------------------------
\subsection{Lesson 3: Check that the adaptation term does what it says}
\label{sec:L3}

\paragraph{Lesson.} An adaptation loss can be numerically inert, act
through a side channel, or have the wrong sign, and the training curve
looks normal in each case.

\paragraph{What happened.} Four failures, each found only by a targeted
check. \todo{(a) Inert: the DANN adversarial gradient was
$\sim 5\times10^{-4}$ of the task gradient; the MMD term was
$\sim$0.2\,\% of the task loss. (b) Pooled away: alignment of
global-average-pooled bottleneck features did nothing; spatially resolved
alignment traded source for target risk. (c) Side channel: every method
with a target forward pass mixes target statistics into the BatchNorm
running estimates. Replaying that traffic alone on an untouched model
reproduces the source degradation in all five pairs tested, and
re-estimating on source removes it, so the weights are barely changed.
AdaBN, which does this on purpose, harmed most directions. (d) Sign
error: our FDA implementation flipped the target mean.}
\ev{audit\_2026-10-29.md A1, A2; GAP ablation results/gap\_ablation/summary/
(read at the 1e-5 floor only, A9); AFM BN check afm\_bn\_check\_seed42/
(PLAN\_COSMO A1, 2026-11-09); overview \S5.2, \S5.3, \S9.1--9.2}

\paragraph{Recommendation.} \todo{Log the ratio of adaptation to task
gradient; evaluate source validation with eval-mode statistics; run a
target-forward-only control at zero alignment weight; unit-test input
transforms on a known shift.}

%---------------------------------------------------------------------
\subsection{Lesson 4: Evaluation can manufacture recovery}\label{sec:L4}

\paragraph{Lesson.} Without target labels for model selection, and with
skewed, zero-inflated targets, the evaluation choices decide more of the
reported recovery than the method does.

\paragraph{What happened.} \todo{(a) The log floor: at $10^{-5}$ the
into-Europe ``gains'' were the model leaving the all-zero predictor, which
both source-only models into Europe essentially are. (b) Physical units:
a DANN model with pixels of $10^{6}$\,mm\,h$^{-1}$ scored well in log
space. (c) References: climatology beats source-only into Europe, so
``better than source-only'' is a weak bar. (d) Units across sources:
source-normalised errors are not comparable across rows. (e) Seeds: the
method ranking changes with the seed in every direction.}
\ev{P3, P14, D1; DANN gr0.5 seed 42; afm\_ensemble\_check summary.csv;
overview \S3, \S5.1--5.2, \S9.3; P12 (Kendall's W) not done}

\paragraph{Recommendation.} \todo{Report zero, climatology and
input-derived (upsampled ERA5 tp) references; a physical-unit column
beside each transformed score; at least three seeds and a rank
concordance; selection rules that do not use target labels.}

% Table candidate: "evaluation traps": trap / how it showed / the check
% that catches it.

%---------------------------------------------------------------------
\subsection{Lesson 5: A probabilistic objective does not change the
diagnosis}\label{sec:L5}

\paragraph{Lesson.} Changing the objective from MSE to a generative model
scored with the CRPS leaves the shape of the transfer gaps intact, but
changes what counts as a gain: spread alone can be rewarded against a
collapsed baseline.

\paragraph{What happened.} \todo{AFM gaps have the U-Net's shape; tropical
sources into Europe collapse the ensemble; FDA into Europe is the only
cell in the project with a partial gain that survives mm space and an
information-free dressed reference, in two seeds. The input-conditioned
ERA5-tp reference decides whether it stays. The oracle into Europe is
close to in-domain for AFM, unlike the U-Net.}
\ev{afm\_seed42/, afm\_seed43/, afm\_seed42\_oracle/;
results/seed42/afm\_ensemble\_check/europe\_west\{,\_tpref\}/ (job 20357462);
overview \S5.6, \S9.3}

\paragraph{Recommendation.} \todo{Score generative models against a
dressed source-only ensemble and an input-conditioned reference; read
spread--skill before crediting a CRPS gain.}

%%%%%%%%%%%%%%%%%%%%%%%%%%%%%%%%%%%%%%%%%%%%%%%%%%%%%%%%%%%%%%%%%%%%%
\section{Discussion}\label{sec:discussion}

\paragraph{When should UDA be expected to help in downscaling?}
\todo{Return to the thesis: shared versus non-shared conditional by
design; where COSMO lands; what that means for applying RainShift-type
models in data-sparse regions. Contrast with structural-realism claims of
generalisation \citep{glawion2025}.}

\paragraph{A checklist.} Table~\ref{tab:checklist} collects the checks
that would have saved us most of the compute spent on this study.

\begin{table}[t]
\caption{Checks before trusting a domain-adaptation result for
downscaling.}\label{tab:checklist}
\begin{center}
\begin{tabularx}{\linewidth}{lX}
\topline
Before training   & Test $P(Y\mid X)$ across regions against a
dependence-preserving null; report common support. (L1) \\
                  & Train the oracle and a few-shot curve. (L2) \\
During training   & Log the adaptation-to-task gradient ratio; eval-mode
source validation; target-forward control at zero alignment weight. (L3) \\
Evaluation        & Trivial references (zero, climatology, upsampled
input); physical-unit bias, wet fraction and maximum; $\geq 3$ seeds and
rank concordance. (L4) \\
Generative models & Dressed and input-conditioned references;
spread--skill. (L5) \\
\botline
\end{tabularx}
\end{center}
\end{table}

\paragraph{Limitations.} \todo{Three RainShift regions and five COSMO
regions; seeds vary initialisation only; two hardware/software stacks
split along the method comparison; U-Net learning rate untuned
(conservative for lesson~5); two binned inputs in lesson~1; every U-Net
is too dry, in-domain included.} \ev{overview \S7; CLAUDE.md traps}

%%%%%%%%%%%%%%%%%%%%%%%%%%%%%%%%%%%%%%%%%%%%%%%%%%%%%%%%%%%%%%%%%%%%%
\section{Conclusions}\label{sec:conclusions}

\todo{One paragraph: the thesis, the five lessons in one sentence each,
and what a practitioner with a data-sparse target should do instead.}

\clearpage
%%%%%%%%%%%%%%%%%%%%%%%%%%%%%%%%%%%%%%%%%%%%%%%%%%%%%%%%%%%%%%%%%%%%%
% ACKNOWLEDGMENTS
%%%%%%%%%%%%%%%%%%%%%%%%%%%%%%%%%%%%%%%%%%%%%%%%%%%%%%%%%%%%%%%%%%%%%
\acknowledgments
\todo{Funding with grant numbers; UNIL Curnagl HPC; the RainShift authors
for the data; the COSMO-CLM simulation authors; OODiagnostics
collaborators.}

%%%%%%%%%%%%%%%%%%%%%%%%%%%%%%%%%%%%%%%%%%%%%%%%%%%%%%%%%%%%%%%%%%%%%
% DATA AND SOFTWARE AVAILABILITY STATEMENT
%%%%%%%%%%%%%%%%%%%%%%%%%%%%%%%%%%%%%%%%%%%%%%%%%%%%%%%%%%%%%%%%%%%%%
\availstatement
\todo{RainShift data source and licence; COSMO-CLM data access (Cui et
al. 2024), and whether the derived 10$\times$ stores can be shared; code
repository and commit; trained weights if released.}

%%%%%%%%%%%%%%%%%%%%%%%%%%%%%%%%%%%%%%%%%%%%%%%%%%%%%%%%%%%%%%%%%%%%%
% APPENDIXES (AMS: \appendix[X], \appendixtitle, then \subsection)
%%%%%%%%%%%%%%%%%%%%%%%%%%%%%%%%%%%%%%%%%%%%%%%%%%%%%%%%%%%%%%%%%%%%%

\appendix[A]
\appendixtitle{Implementation details}
\subsection*{Models, training and splits}
\todo{Architectures, hyperparameters, validation split (strided, purge
gap), compute and the two software stacks.}

\appendix[B]
\appendixtitle{Audit log and dropped methods}
\subsection*{What was found and what it changed}
\todo{Each bug or inert configuration found, how it was found, and what
it changed: FDA sign, DANN gradient ratio, AdaBN scoring the unadapted
checkpoint, CORAL and multi-scale MMD, BatchNorm mixing. This appendix is
the evidence base for lesson~3.}
\ev{report/audit\_2026-10-29.md}

\appendix[C]
\appendixtitle{Per-direction results}
\subsection*{Full tables}
\todo{Full transfer matrices, oracle and method tables at the 0.1\,mm
floor with mm columns, both architectures, all seeds.}

%%%%%%%%%%%%%%%%%%%%%%%%%%%%%%%%%%%%%%%%%%%%%%%%%%%%%%%%%%%%%%%%%%%%%
% REFERENCES
%%%%%%%%%%%%%%%%%%%%%%%%%%%%%%%%%%%%%%%%%%%%%%%%%%%%%%%%%%%%%%%%%%%%%
\clearpage

\bibliographystyle{ametsocV6}
\bibliography{references}

\end{document}
